Image quality is a fundamental requirement for reliable biometric recognition and secure identity systems.
Yet, existing evaluation methods often fail against distortions that remain invisible to the human eye but critically affect recognition, such as those introduced by steganography.
They also reproduce demographic biases, leading to uneven treatment across age, gender, and skin tone.
This dissertation introduces a pipeline for creating new image quality metrics through a progressive path from full-reference fusion to no-reference prediction.
The first stage combines complementary full-reference measures into a single fused metric that reflects human perception more faithfully than any individual method.
This produces metrics such as pMOS-SSIM+, which show superior alignment with subjective opinion scores across benchmark datasets.
The second stage uses the fused metric as a proxy ground truth to supervise no-reference predictors.
This step yields task-specific models, such as pMOS-StegaFace, tailored to steganographically distorted facial images, and pMOS-NoisyTID, which generalizes to diverse content and distortion types.
Results confirm that the fused metrics consistently outperform established baselines, and that the no-reference models trained from them detect subtle degradations while maintaining strong correlation with human judgments.
The pipeline also integrates fairness-aware debiasing of subjective scores, ensuring consistent performance across demographic groups.
By uniting perceptual validity, biometric utility, and operational security, this work establishes a flexible methodology: one may adopt the fused metric as a new full-reference standard, or extend it into a no-reference model.


% Objective image quality assessment often diverges from human evaluation on subtle, security-relevant degradations in face images, notably steganographic perturbations that are hard to see, yet harmful for biometric pipelines. This thesis proposes a Full-to-No Reference Facial Image Quality Assessment framework that learns a perceptually aligned no reference metric from a small set of human ratings. First, we run a single stimulus subjective study on a labeled subset and remove demographic effects from Mean Opinion Score via ordinary least squares residualization validated with ANOVA and MANOVA. Second, we train a fusion regressor on the most relevant full reference metrics to approximate MOS and generate pseudo MOS for the unlabeled majority. Third, we train an image based regressor on MOS plus pseudo MOS to obtain the final no reference FIQA score. Experiments on ICAO compliant FRLL faces degraded by four steganographic methods at nine intensities show that the best fusion reaches PLCC around 0.89 with MOS on the labeled subset, and the learned no reference metric attains PLCC around 0.93 and SRCC around 0.93 on disjoint identities, outperforming NIQE, PIQE, BRISQUE and FIQA baselines such as FaceQnet, while SER FIQ and MagFace exhibit negative correlation with human opinion. The approach needs human labels for roughly ten percent of the data, yields a reproducible and bias aware pipeline, and offers a practical path to perceptually grounded FIQA when authentic distortions are subtle and annotations are scarce.


% Face Image Quality Assessment (FIQA) is meant to predict whether a face image will work well in recognition systems. In practice, widely used FIQA scores often disagree with what people perceive as “good” or “poor” quality—especially for images that look fine at a glance but contain subtle artifacts (e.g., mild blur, compression, color shifts) that still affect recognition. This thesis targets that misalignment.

% We introduce a “full-to-no-reference” pipeline that learns a perceptually aligned FIQA score from a small set of human ratings. First, we run a single-stimulus subjective study and debias the resulting mean-opinion scores (MOS) for demographic effects via ordinary least squares residualization, validated with ANOVA/MANOVA. Second, we train a lightweight fusion regressor on a compact set of full-reference metrics to approximate MOS, then use it to generate pseudo-MOS for the much larger unlabeled set. Third, we train a no-reference model on MOS + pseudo-MOS to produce the final FIQA score—one that requires no reference image at test time but tracks human perception.

% On ICAO-compliant FRLL faces degraded across nine intensity levels, the best full-reference fusion reaches ~0.89 Pearson correlation with MOS on the rated subset. The learned no-reference FIQA generalizes to disjoint identities with ~0.93 Pearson and ~0.93 Spearman correlation, outperforming common no-reference baselines (NIQE, PIQE, BRISQUE) and FIQA baselines such as FaceQnet; SER-FIQ and MagFace show negative correlation with human opinion in this setting. The approach uses human labels for only ~10% of the data, yields a reproducible and bias-aware pipeline, and provides a practical path to FIQA that better reflects human perceptual judgments.